% Generated draft: requires mathematical and source review.
\documentclass[11pt]{article}
\usepackage[margin=25mm]{geometry}
\usepackage{fontspec}
\setmainfont{Latin Modern Roman}
\usepackage{amsmath,amssymb,mathtools,unicode-math}
\setmathfont{Latin Modern Math}
\usepackage{enumitem,needspace,xurl,hyperref}
\setlength{\parindent}{0pt}
\setlength{\parskip}{5pt}
\setlength{\emergencystretch}{3em}
\newcommand{\N}{\mathbb{N}}
\newcommand{\Z}{\mathbb{Z}}
\newcommand{\Q}{\mathbb{Q}}
\newcommand{\R}{\mathbb{R}}
\newcommand{\C}{\mathbb{C}}
\newcommand{\F}{\mathbb{F}}
\newcommand{\A}{\mathbb{A}}
\newcommand{\PP}{\mathbb{P}}
\newcommand{\E}{\mathbb{E}}
\newcommand{\eps}{\varepsilon}
\newcommand{\dd}{\,\mathrm{d}}
\begin{document}
\section*{UnsolvedMath ID: 30005646}
\textbf{Source problem code:} catalogue entry
\subsection{Definitions and mathematical statement}
Let a global field mean a finite extension of \(\Q\) or of \(\F_q(t)\), where \(\F_q\) is a finite field. For a global field \(K\), let \(v\) range over its places, let \(K_v\) be its completions, and let \(\A_K\) be its adele ring. For a variety \(X\) over \(K\), write \(X(\A_K)\) for its set of adelic points. Following the supplied record, write
\[
X(\A_K)^h\subseteq X(\A_K)^{h,2}\subseteq X(\A_K)
\]
for the subsets surviving, respectively, the full \(\acute{e}\)tale-homotopy obstruction and its two-stage truncation.

The requested existence statement is
\[
\exists\,K\text{ a global field}\;\exists\,X/K\text{ a variety}\;\exists\,a\in X(\A_K):\qquad
 a\in X(\A_K)^{h,2}\ \text{and}\ a\notin X(\A_K)^h.
\]
Equivalently, find \(K\) and \(X\) for which
\[
X(\A_K)^h\subsetneq X(\A_K)^{h,2}.
\]
The supplied excerpt does not define the obstruction subsets or its truncation convention, and does not specify whether the varieties must be smooth, proper, geometrically integral, or geometrically connected. Consequently, the notation above retains the source's intended obstruction constructions; a fully self-contained formulation and the admissible class of varieties require verification against the report.

The original excerpt subsequently begins a separate discussion of higher dimensional global fields. It describes \(K=k(C)\), for a curve \(C/k\), as having cohomological dimension one greater than that of \(k\). For \(k/\Q_p\) finite, it mentions Poitou--Tate duality using the closed points of \(C\) as places, and claims an identification of \(X(\A_K)^{Zh}\) with an analogue of the Brauer--Manin obstruction involving \(H^3(X;\Q/\Z(2))\). It also writes \(\operatorname{Br}X=H^2(X;\Q/\Z(1))\). These are supplied background claims with unspecified hypotheses and conventions, not additional requirements of the existence statement or verified identities. In particular, the meaning of the superscript \(Zh\) and the cohomological identifications remain unresolved.
\subsection{Short English statement}
Find a variety over a global field with an adelic point that survives the two-stage \(\acute{e}\)tale-homotopy obstruction but fails the full obstruction. The supplied record labels this problem open; its precise obstruction conventions and geometric hypotheses require human mathematical review.
\subsection{Sources}
\begin{itemize}
\item Supplied local UnsolvedMath registry record, permanent ID \(30005646\), source problem code \texttt{catalogue entry}, in the Oberwolfach Reports Open Problems collection. The local registry paths designated by the project are \texttt{data/unsolvedmath/problems.json} and the supplementary \texttt{catalog/problems.json}; no file lookup was performed. No individual proposer is supplied.
\item Supplied source citation: \emph{MFO--RIMS Tandem Workshop: Arithmetic Homotopy and Galois Theory (2024)}, Oberwolfach Reports, pp.\ 2377--2488, DOI \texttt{10.4171/owr/2023/42}. \url{https://doi.org/10.4171/owr/2023/42}. The record gives proposed year \(2023\); the citation's \(2024\) label and the \(2023\) provenance require reconciliation.
\item Supplied literature source, labelled ``Schlank, etale homotopy obstructions of arithmetic spheres'': \url{https://arxiv.org/abs/1902.03404}. This attribution and the link's bibliographic correspondence have not been independently verified.
\item Supplied literature source, labelled ``OWR report'': \url{https://doi.org/10.4171/owr/2023/42}.
\item Literature assessment dated exactly \texttt{2026-08-22}: supplied dated curation, not independently verified. It labels the problem open, reports that the workshop poses strictness as a natural question and cites only progress in work in progress, and says that no demonstrable general example was found. The background describes this as dated literature triage and requests independent verification of time-sensitive status and cited claims.
\item The registry assigns \texttt{statement\_status: corrected\_verified} and supplies a claim that the clean statement was checked against the official report and preserves its hypotheses, notation, and target. These are supplied verification metadata, not verification performed for this draft.
\end{itemize}
\end{document}
